\section*{Chapitre \ref{ch:m3:emissions-and-environmental-markets} --- Émissions et marchés environnementaux}

\begin{solution}{exo:m3:emissions-and-environmental-markets:1}
Il restitue 1\,000 quotas, au plus tard le 30 septembre de l'année suivant les émissions~; les 100 achetés
couvrent le déficit.
\end{solution}

\begin{solution}{exo:m3:emissions-and-environmental-markets:2}
À 1\,200 millions, il retient 24~\%, soit 288 millions~; à 900 millions, l'excédent au-delà de 833 millions, 67
millions~; à 350 millions, il libère 100 millions.
\end{solution}

\begin{solution}{exo:m3:emissions-and-environmental-markets:3}
Avec une taxe, le prix reste fixe et les émissions baissent avec la demande~; avec un plafond, la quantité reste au niveau du plafond
et le prix baisse, laissant un excédent qui peut déprimer les prix pendant des années, à moins qu'un mécanisme comme la
réserve ne réduise l'offre future.
\end{solution}

\begin{solution}{exo:m3:emissions-and-environmental-markets:4}
\qty{13.11}{EUR/t} avec le gaz à \qty{20}{}~; \qty{60.47}{EUR/t} avec le charbon à 38~\% et le gaz à
\qty{35}{}~: la centrale au charbon, moins efficace, émet davantage par MWh, si bien qu'un prix du carbone plus bas suffit.
\end{solution}

\begin{solution}{exo:m3:emissions-and-environmental-markets:5}
$0.202/0.55 = 0.367$~t par MWh~; une hausse de 10~euros ajoute \qty{3.67}{EUR/MWh}.
\end{solution}

\begin{solution}{exo:m3:emissions-and-environmental-markets:6}
Une \omterm{def:m3:emissions-and-environmental-markets:go}{garantie d'origine} transfère une affirmation sur la façon dont un MWh a été produit~; le MWh renouvelable a été produit
de toute façon, et aucun plafond ne contraint l'acheteur. Un quota restitué est annulé dans un système de plafonnement~: le total
pouvant être émis est fixé, de sorte que chaque tonne compensée est une tonne de moins disponible pour quelqu'un d'autre.
\end{solution}

\begin{solution}{exo:m3:emissions-and-environmental-markets:7}
Entre 833 et 1\,096 millions, le rapport $(T - 833)/T$ croît avec $T$, pour atteindre 23.996~\% à 1\,096
millions~; au-delà, le prélèvement vaut 24~\% de $T$. À 1\,096 millions, les deux morceaux donnent 263.00 et 263.04
millions~: la fonction est continue à 0.04 million près.
\end{solution}

\begin{solution}{exo:m3:emissions-and-environmental-markets:8}
Un prix bas peut signifier que le plafond est respecté à moindre coût (réduction peu coûteuse, demande faible), et non que
le plafond ne contraint pas~: les émissions ne peuvent dépasser le plafond quel que soit le prix. Le prix mesure le coût marginal de
rester sous le plafond, et l'excédent qu'il laisse est reporté aux années suivantes, où le plafond sera plus serré.
\end{solution}

\begin{solution}{pb:m3:emissions-and-environmental-markets:1}
\textbf{1.} Charbon $0.341/0.40 = 0.8525$~t par MWh~; gaz $0.202/0.55 = 0.3673$~t. \textbf{2.} $6 \times
0.8525 + 4 \times 0.3673 = 6.584$ millions de tonnes. \textbf{3.} Environ 460.9 millions d'euros. \textbf{4.} Au plus tard le
30 septembre de l'année suivante. \textbf{5.} Pour figer le coût lorsqu'elle vend à terme son électricité, et parce
qu'attendre l'expose au prix pendant jusqu'à 21 mois.
\textbf{6.} \qty{69.32}{EUR/t}. \textbf{7.} C'est le gaz, moins cher de \qty{0.33}{EUR/MWh} (89.35 contre 89.68).
\textbf{8.} $(0.8525 - 0.3673) \times 10^6 = 485\,227$ tonnes. \textbf{9.} Au-dessus du \omterm{def:m3:emissions-and-environmental-markets:switch}{prix de bascule},
\qty{69.32}{EUR/t}, en négligeant les autres coûts de chaque centrale. \textbf{10.} Elle monte à \qty{125.53}{EUR/t}~:
le charbon devient la centrale la moins coûteuse à \qty{70}{}.
\textbf{11.} Avec un TNAC de 1\,023\,494\,202, la réserve retient 190\,494\,202 quotas des
mises aux enchères de septembre 2026 à août 2027~: l'offre baisse. \textbf{12.} Une demande plus faible fait baisser le prix
et fait monter le TNAC, ce qui accroît les prélèvements futurs et resserre l'offre plus tard. \textbf{13.} Le coût de la centrale marginale
comprend $0.8525 \times P_{\mathrm{CO_2}}$ par MWh, ce que la mise aux enchères répercute dans le prix.
\textbf{14.} Les investisseurs financiers et les fonds, les spéculateurs et les intermédiaires qui vendent à terme aux
acheteurs soumis à l'obligation de restitution~; les banques qui couvrent les fournisseurs d'énergie. \textbf{15.} Rien directement pour ses ventes
d'électricité intérieures, mais l'électricité importée et les biens à forte intensité énergétique supportent désormais un coût carbone, ce qui modifie la
concurrence. \textbf{16.} Sa marge est un clean spread~: le prix de l'électricité, celui du combustible et celui du carbone évoluent ensemble
dans l'\omterm{def:m3:power-markets-design:merit}{ordre de préséance économique} (merit order), et couvrir une seule jambe laisse le spread ouvert. \textbf{17.} Que la
centrale ne dégage guère de marge au-delà du combustible et du carbone et fonctionnera moins à mesure que le plafond se resserre~: fermeture ou conversion. \textbf{18.}
Quand le gaz est bon marché, la substitution du gaz au charbon est la réduction d'émissions la moins coûteuse à grande échelle, si bien que le \omterm{def:m3:emissions-and-environmental-markets:switch}{prix de bascule} joue le rôle de
plafond que le prix du quota dépasse rarement tant que la capacité de substitution subsiste. \textbf{19.}
\emph{Résultat nommé~:} un \omterm{def:m3:emissions-and-environmental-markets:switch}{prix de bascule} de \qty{69.32}{EUR/t}, et 6.584 millions de quotas (environ
460.9 millions d'euros à \qty{70}{EUR/t}) pour le plan. \textbf{20.} Le coût de la tonne marginale
de réduction nécessaire pour rester sous le plafond, aujourd'hui et en espérance (les quotas peuvent être reportés).
\end{solution}

\begin{solution}{iq:m3:emissions-and-environmental-markets:1}
Une taxe fixe le prix par tonne et laisse les émissions s'ajuster~; le \omterm{def:m3:emissions-and-environmental-markets:cap}{système de plafonnement et d'échange} fixe la quantité et laisse le
prix s'ajuster par les échanges. Face à l'incertitude sur les coûts de réduction, l'un fixe le résultat, l'autre le
coût.
\iqlookfor{instrument prix contre instrument quantité, et qui supporte l'incertitude.}
\end{solution}

\begin{solution}{iq:m3:emissions-and-environmental-markets:2}
Le prix de l'électricité moins les coûts du gaz et du carbone pour un MWh d'une centrale à gaz. Négatif lorsque l'électricité est
moins chère que le combustible de la centrale plus les quotas~: aux heures fixées par les énergies renouvelables ou par des centrales moins coûteuses.
\iqlookfor{la formule avec le rendement et le facteur d'émission.}
\end{solution}

\begin{solution}{iq:m3:emissions-and-environmental-markets:3}
Elle fait baisser l'offre future aux enchères à mesure que l'excédent augmente~: un TNAC élevé retire des quotas, un TNAC faible
les libère. L'offre devient une fonction des propres avoirs du marché, ce qui amortit les
excédents persistants et soutient les prix après les chocs de demande.
\iqlookfor{offre endogène.}
\end{solution}

\begin{solution}{iq:m3:emissions-and-environmental-markets:4}
Un gaz plus cher relève le \omterm{def:m3:emissions-and-environmental-markets:switch}{prix de bascule}, si bien que davantage de charbon tourne et que l'électricité émet davantage par MWh~: la demande de
quotas augmente, et les prix ont tendance à monter~; en sens inverse, des prix de l'énergie élevés réduisent la production industrielle et
la demande. L'effet net dépend de celui qui domine.
\iqlookfor{substitution de combustible comme canal de demande, et effet d'activité compensateur.}
\end{solution}

\begin{solution}{iq:m3:emissions-and-environmental-markets:5}
Risque de qualité (crédits dont la réduction n'a pas eu lieu ou aurait eu lieu de toute façon), risque de permanence
(forêts qui brûlent), risque de réputation et risque juridique (allégations jugées trompeuses), et risque de prix si l'on doit
remplacer des crédits invalidés.
\iqlookfor{additionnalité, permanence, allégations.}
\end{solution}

\begin{solution}{iq:m3:emissions-and-environmental-markets:6}
Un modèle de dispatching~: demande horaire, scénarios de production renouvelable, le parc de centrales avec leurs rendements et leurs facteurs
d'émission, scénarios de prix du combustible et du carbone, et interconnexion~; appliquer l'\omterm{def:m3:power-markets-design:merit}{ordre de préséance économique} heure par heure et sommer
les émissions. Données~: registres du parc, dispatching historique, météorologie. Calibrer sur les émissions vérifiées passées.
\iqlookfor{un modèle de dispatching à \omterm{def:m3:power-markets-design:merit}{ordre de préséance économique} et scénarios, calibré sur les émissions vérifiées.}
\end{solution}
